\subsection{The Selmer datum for a main critical line}\label{moment:line}

Fix a main pair $(\nu,z)$ at stage $1\leq j\leq m$, and put $t_0=2z-1$.
Write $p_0,p_1$ for the two distinct level primes and $f=[E:\mathbb Q]/2$.
The tame character of $\nu$ at each chosen place is the $j$th power of
the prescribed residue character of order $q_2=(q_0d)^{s+6}$.
There is one chosen place from each conjugate pair above each $p_h$,
and $p_h\equiv 1\pmod {q_2^2}$. These are the level and main-pair
conditions of the construction. In particular $p_h$ splits completely in $E$,
whose conductor divides $q_1d\mid q_2$. Enlarge the coefficient field
at $\ell$ to contain the character values and all labelled embeddings.
Let $\mathcal O$ be its integer ring, $\pi$ a uniformizer, and $T$ the
stable line of \S\ref{rankone:main}, with character $\tau$.
Its weights at a pair of complex embeddings are $z-1-c_v$ and $z-1+c_v$,
and $\tau^c\tau=\epsilon_\ell^{1-t_0}$.
The only places at which the line can ramify are those above
$\ell p_0p_1$.

\paragraph{Local rank at $\ell$.}
For $v\mid\ell$, let $Z_v$ be the continuous $T$-valued cocycles whose
rational cohomology class belongs to the Bloch--Kato group
$H^1_f(E_v,V)$, $V=T[1/\ell]$. This specifies the integral condition
in G.2 by imposing, at every torsion level, the homogeneous evaluation
relations valid on $Z_v$. Coboundaries lie in $Z_v$. The module is
torsion-free and finite over $\mathcal O$: the coboundaries are an
image of $T$ and $H^1(E_v,T)$ is finitely generated by local Galois
cohomology. Rational cocycles on the compact decomposition group
have bounded image, so tensoring $Z_v$ with the coefficient field
recovers all rational cocycles with finite class. The character is
potentially crystalline, hence de Rham. The Bloch--Kato dimension formula
(\href{https://virtualmath1.stanford.edu/~conrad/BSDseminar/refs/BKTamagawa.pdf}{\emph{L-functions and Tamagawa numbers of motives}, Corollary 3.8.4}~\cite{bk})
gives $\dim H^1_f=\dim H^0+\dim(D_{\rm dR}/\mathrm{Fil}^0)$.
Adding the coboundary dimension $1-\dim H^0$ gives
\[
  \operatorname{rank}_{\mathcal O} Z_v
  =1+\#\{\sigma\mid v : \operatorname{HT}_{\sigma}(V)<0\}.
\]
Since the two weights at a conjugate pair sum to $2z-2\geq0$, at most
one is negative. Summing yields the bound $\#\{v\mid\ell\}+f$
required in G.2.

\paragraph{An inertia scalar and the old local terms.}
Write $Q=q_2$. At a chosen $w\nmid\ell$ over $p_h$ the value $\xi$
of $\tau$ on a tame generator has order $Q/\gcd(Q,j)$ by the
rank-one formula under local reciprocity. Inversion of a generator or
character preserves that order. Choose $a\in\{2,5\}$ different from
$\ell$. Since $a\mid q_0$, we have $a^s\mid Q$ and
$a^s\ge 2^s=m^{22}>2m$ on the constructed scales. If the residual
order of $\xi$ divided $2$, then $Q\mid 2j\ell^b$ for some $b$,
which would imply $a^s\mid 2j$. Thus $\xi^2-1$ and $\xi-1$ are
units. Inertia at $w$ has trivial cyclotomic character, so
$\tau(\sigma^{-c})=\tau(\sigma)=\xi$ for the chosen generator $\sigma$;
here $\sigma^{-c}=(c\sigma c^{-1})^{-1}$. This also handles the
conjugate place. At least one of $p_0,p_1$ differs from $\ell$,
so such a $\sigma$ always exists. The matrix algebra condition in
rank one holds since the identity is in the image.

For $W=T/\pi^u T$ the local Tate dual has no invariants at any such
$w$: inertia acts on it by $\xi^{-1}$ on a generator with $\xi-1$
a unit. Local duality and the Euler characteristic at $w\nmid\ell$
then give $\operatorname{length} Z^1(E_w,W)=\operatorname{length}W$.
Thus the restricted old local cocycles of G.2 have zero positive excess.
The two determinant losses for releasing the first auxiliary prime also
vanish if its Frobenius and inverse conjugate have the prescribed
residual inertia values. This is an open congruence condition in G.2.

\paragraph{The global term.}
The conjugate Tate dual of $W$ has character
$\epsilon_\ell(\tau^c)^{-1}=\tau\epsilon_\ell^{t_0}$.
As an $\mathcal O$-module the Tate dual is free of rank one over
$\mathcal O/\pi^u$: the trace-dual fractional ideal of $\mathcal O$
is principal. Choose $p_h\ne\ell$. For an arithmetic Frobenius lift
at a degree-one place over $p_h$ the cyclotomic value is $p_h$;
ramification of $T$ there does not affect the displayed character
quotient. Every equivariant map from $W$ to its conjugate Tate dual
therefore has image killed by $p_h^{t_0}-1$. Its value on a basis vector
determines the map. Uniformly in $u$, the length of the space of
equivariant maps, divided by $v_\pi(\ell)$, is at most
$v_\ell(p_h^{t_0}-1)$. The group-order budget covers this integer:
$t_0=2z-1\le 2r-1\le r t$ for $t=5r$, and the factor
$p_h^{t_0}-1$ occurs in the order of the chosen
$\mathrm{GL}_{rt}(\mathbb F_{p_h})$.
We have thus verified the rank, scalar, and length hypotheses in G.2
for the explicit untwisted line with the stated local conditions.
The residue-field size condition in G.1 can be met by an unramified
coefficient extension at each fixed torsion level (its ramification
index is unchanged).

\paragraph{Split Hecke action on the lift of the constant.}
Put $N=2z$ and take a split $w\nmid\ell$ away from the level
and the conductor of a tame twist $\eta$. Write $q=\#k_w$.
Use the split Schr\"odinger model with local vacuum
$\mathbf1_{\mathcal O_w^{N\times1}}$. For a splitting character $a$
of rank $d$ its Frobenius factor is
\[
 c_d(a)=q^{d/2}a(\varpi_w/\bar\varpi_w)\upsilon(\varpi_w)^{-d}
\]
with $\varpi_w$ interpreted as a local uniformizer idele, and
$\upsilon$ the unitary splitting Hecke character. Write
$\alpha=c_1(a_W)$ and $\beta=c_N(a_V)$ for the source and target
factors of the pair. Conjugation of the local variable on the source
side corresponds to the adjoint Hecke action.

Normalize the Hecke polynomial in rank $d$ by
$P_d(X)=\sum_i(-X)^i q^{\binom{i}{2}}T_i$ for the unaveraged
double-coset sums. The cosets in $T_i$ are indexed by
$i$-planes $U\subset k_w^d$: one can take a representative
$\gamma_U$ integral with elementary divisors $\varpi_w$ on $i$ axes
and $1$ elsewhere and with reduction kernel $U$.
For a matrix $y$ of width $d$, $y\gamma_U^{-1}$ is integral exactly
when $y$ is integral and $\bar y$ kills $U$. Indeed, in adapted
integral bases the inverse divides the coordinates of $U$ by
$\varpi_w$. If $y\gamma_U^{-1}$ is integral, multiplication by
$\gamma_U$ gives integrality of $y$.

On the target vacuum the transpose coset therefore gives a multiplier
$\alpha^i$ times the indicator of integrality of $y\gamma_U^{-1}$,
where $y$ is the transpose of the local column. On the dual source
vacuum the inverse coset gives the same expression with $y$ the
column and multiplier $\beta^i$. These are the split Weil
change-of-variables formulas: the determinant has valuation $i$,
its unramified splitting-character value is the $i$th power of the
uniformizer value, and the inverse square-root Jacobian contributes
$q^{di/2}$ with $d$ the rank of the opposite space. Thus for integral
$y$ of residual rank $a$ and width $d$, the resulting weighted sum,
with $Z$ incorporating the character and Jacobian, is
\[
 \sum_{i=0}^{d-a}{d-a\brack i}_q(-Z)^i q^{\binom{i}{2}}
 = (Z;q)_{d-a},\qquad (Z;q)_h=\prod_{i=0}^{h-1}(1-q^iZ).
\]
Here the Gaussian binomial counts the $i$-planes in $\ker\bar y$;
the finite $q$-binomial identity follows by expanding the product.
For a nonintegral matrix every term vanishes. The column and its
transpose have the same rank $a\in\{0,1\}$. Consequently
\[
 (Z;q)_{N-a}=(Z;q)_{t_0}(q^{t_0}Z;q)_{1-a}.
\]
On the target take $Z=\alpha X$. The polynomial argument on the
source is then $q^{t_0}\alpha X/\beta$.

For the global passage, suppose the chosen section of the theta
kernel has this vacuum component at $w$, and that the global lift
pairs the source form and the kernel equivariantly under their finite
translations. In the usual theta integral this equivariance follows
by change of integration variable with invariant measure, provided
the kernel and the chosen section have the specified transformation
law. It moves a target Hecke sum onto the vacuum and then moves the
adjoint source sum onto the source form. The constant source on $U(1)$
has Hecke polynomial $1-X$, since the sole rank-one uniformizer
translation fixes it. Hence its lift has reciprocal Hecke polynomial
\[
 (\alpha X;q)_{t_0}\bigl(1-q^{t_0}\alpha\beta^{-1}X\bigr).
\]
Multiplying the global lift by the source level index, as in the
normalization of the manuscript's character lift, preserves this
identity. A nonzero lift is required to read its eigenvalues.

The negative linear coefficient, after the target normalizing twist
by $\alpha^{-1}$ at geometric Frobenius, is
$q^{t_0}/\beta+\sum_{a=0}^{t_0-1}q^a$.
To check the normalization, write $u=\nu\eta$. In even rank the
splitting convention is $a_V=u\,\upsilon_{U(1)}^{z}$, with
$\upsilon_{U(1)}(x/\bar x)=\upsilon(x)^2$. Thus
$\beta=q^{z}u(\varpi_w/\bar\varpi_w)$.
Consequently the first term is $q^{z-1}u(\varpi_w/\bar\varpi_w)^{-1}$,
the uniformizer value of $u(x/\bar x)^{-1}|x|^{1-z}$ in \S\ref{rankone:main}.
The other terms are the values of $\epsilon_\ell^{-a}$ on geometric
Frobenius. If a continuous representation attached to the target
has these Hecke polynomials, on the density-one set of degree-one
unramified primes of $E$ this gives the normalized trace
\[
 \eta^{-1}_{\ell}\tau+\sum_{a=0}^{2z-2}\epsilon_\ell^{-a},
\]
where $\eta_\ell$ is the avatar of $x\mapsto\eta(x/\bar x)$.
Degree-one primes away from ramification are split over $E^+$;
the remaining primes of absolute residue degree at least two have
convergent reciprocal-norm sum. Chebotarev~\cite[Theorem~8.31]{milneANT} and continuity extend this
trace equality to $G_E$, and characteristic-zero character theory
identifies semisimplifications.

The local computation uses the displayed split action. For its
application to the manuscript's lift one must identify the actual
unramified section and the equivariant global pairing with these
normalizations, and supply the attached target representation.
The integral extension classes and generator-wise clearing operator
require separate first-order congruences. For the higher-rank branch
one also needs ramified label compatibility and the full integral
matrix span for a lattice in the main summand.

\paragraph{Source correspondence.}
At revision \texttt{72fead1}, the calculations can be compared with
the
\ifdefined\href\href{https://github.com/lsernau/grh-lean/blob/72fead102f663d5c7a407964a13dc0cf9d977207/grh/src/New/Congruences/Selmer/MainSummands.lean\#L294}{Selmer datum for a main summand}\else Selmer datum for a main summand\fi,
the
\ifdefined\href\href{https://github.com/lsernau/grh-lean/blob/72fead102f663d5c7a407964a13dc0cf9d977207/grh/src/Lib/Weil/Unramified/HeckeTransfer.lean\#L413}{split vacuum Hecke calculation}\else split vacuum Hecke calculation\fi,
the
\ifdefined\href\href{https://github.com/lsernau/grh-lean/blob/72fead102f663d5c7a407964a13dc0cf9d977207/grh/src/New/OneLift/HeckeValues/CharacterLiftNorm.lean\#L215}{specialization to the constant lift}\else specialization to the constant lift\fi,
and the
\ifdefined\href\href{https://github.com/lsernau/grh-lean/blob/72fead102f663d5c7a407964a13dc0cf9d977207/grh/src/New/Congruences/GaloisRepresentations/RowsFromEigensystems.lean\#L340}{normalized Frobenius trace identity}\else normalized Frobenius trace identity\fi.
